% Abstract – literature review portion only (博士生1)
% To be merged into the full abstract by the lead editor.
% 2–4 sentences: prior knowledge, key studies/data, gap this paper fills.

Existing work on China's food away from home (FAH) or out-of-home food consumption has relied mainly on household surveys and expenditure data, documenting rising FAH shares of total food spending and income-driven growth in urban dining-out, but has not produced nationally or provincially representative estimates of out-of-home consumption \emph{by food commodity} that can be used for macro-level food balance and food security analysis. International sources such as FAO, OECD, and IFPRI report aggregate food use or consumption trends for China yet typically do not distinguish home versus away-from-home consumption by product, and survey-based studies remain difficult to scale to provincial or national coefficients. This paper fills that gap by providing the first machine-learning-based, commodity-level estimates of out-of-home consumption coefficients for China at provincial and national levels, bridging micro survey data and macro prediction scenarios through Copula income-distribution matching, spatial interpolation, and multi-model comparison with uncertainty quantification.
